\documentclass{article}
\usepackage{amsmath,amssymb,mathtools}
\usepackage[margin=1in]{geometry}
\allowdisplaybreaks
\usepackage[UTF8]{ctex}
\begin{document}

\section*{DFRP无限时域折扣马尔可夫决策模型}

\[
\pi^*=\arg\max_{\pi}\mathbb{E}\left\{\left.
\sum_{t=0}^{\infty}\gamma^t C_t\bigl(S_t,X^\pi(S_t)\bigr)
\right|S_0\right\}.
\]
优化过程是在确定性静态策略的基础上进行的：策略$\pi$将每个状态$S_t$映射到一个可行的路由决策上。
\[
X^\pi(S_t)\in\mathcal{X}(S_t)\qquad\text{for every state }S_t.
\]
这里的$S_0$是一个初始状态。
\[
0<\gamma<1.
\]
平稳性是指相对于整个状态而言的，包括周期中的每一天。因此，当仅将政策视为库存组成部分$R_t$的函数时，政策可能会因不同的天数而有所变化。

\subsection*{可行的路由决策}
\begin{equation}
\mathcal{X}(S_t)=\left\{x_t:\begin{aligned}
&\sum_{j\in\mathcal{L}:((i,[t]),(j,[t+1]))\in\mathcal{A}}x_{ijt}^{dl}=R_{it}^{dl}
&&\forall i,d\in\mathcal{L},\ l\in\mathcal{T},\\
&\sum_{l\in\mathcal{T}}\sum_{d\in\mathcal{L}}x_{ijt}^{dl}\leq Q_{ij[t]}
&&\forall ((i,[t]),(j,[t+1]))\in\mathcal{M},\\
&x_{ijt}^{dl}\in\mathbb{Z}_{\geq0}
&&\forall l\in\mathcal{T},\ i,j,d\in\mathcal{L},\\
&x_t\in\mathcal{Y}
\end{aligned}\right\}.
\tag{1}
\end{equation}
固定的集合$\mathcal{Y}$包含了那些符合规定的2-alt负载计划的路由决策。战术计划是固定的，适用于整个运营季节：它规定了每周的拖车移动模式，以及每个地点两个后续终端的选择方案。这些后续终端的选择方案取决于当前的位置和最终目的地，而不是货物的出发地。来源点$\mathcal{Y}$是一个抽象的概念；在MDP公式中，它并没有被赋予具体的代数表达式。

\subsection*{固定的输入、网络以及需求}
时间被划分为不同的运营日。每个运营日的开始是一个决策时刻。运营时间范围是无限的，而拖车移动模式以及随机需求模式则会在每个周期中重复出现。这个周期的长度是$T$天。
\[
\mathcal{T}=\{0,1,2,\ldots,T-1\}.
\]
这种环绕式的时间-空间网络及其节点是
\[
\mathcal{G}=(\mathcal{N},\mathcal{A}),\qquad
n=(i,\tau),\quad i\in\mathcal{L},\ \tau\in\mathcal{T},
\]
其中，$\mathcal{L}$表示终端位置集合。弧线集合由移动弧线$\mathcal{M}$和库存弧线$\mathcal{I}$组成。
\[
\mathcal{A}=\mathcal{M}\cup\mathcal{I}.
\]
运动弧线的形式如下：
\[
((i_1,\tau_1),(i_2,\tau_2)),\qquad i_1\ne i_2,
\]
并表示从$i_1$出发的预定拖车运输在$\tau_1$日至$i_2$日，以及$\tau_2$日的运输。所有运输路径$m\in\mathcal{M}$的行驶时间为一天。在本版本中，将基础时间周期的倍数时间通过向$\mathcal{L}$添加虚拟节点，并将多周期路径分解为单周期路径来处理。因此，货物在决策时刻位于终端，而不是在拖车上运输。从$T-1$日出发的运输路径则循环至该日。
\[
(T-1+1)\bmod T=0.
\]
对于每一个运动弧线：
\[
m=((i_1,\tau_1),(i_2,\tau_2))\in\mathcal{M},
\]
$Q_{i_1i_2\tau_1}$是战术计划所指定的、在预定拖车上能够容纳的容量。库存弧线以打印形式呈现。
\[
((i_1,\tau_1),(i_2,\tau_1+1)),\qquad i_1=i_2,
\]
采用网络的环绕解释方式。它们的持续时间为一天，且无容量限制：在终端处持有托盘始终是一个可行的选择。

没有时间维度的扁平网络是
\[
\mathcal{G}^{flat}=(\mathcal{L},\mathcal{A}^{flat}),
\]
其中，$\mathcal{A}^{flat}$描述了终端之间的连接方式。

该商品由多个托盘组成，这些托盘的共同来源为$o$，目的地为$d$，可用日期为$e$，到期日为$l$。此外，还有$o,d$个终端位置以及$e,l\in\mathcal{T}$个其他信息。其随机数量为
\[
D_{o,d,e,l}\in\mathbb{Z}_{\geq0},
\]
在有限的支持和给定的分布信息下。这是从$o$出发，在$e$日产生的托盘数量，并承诺在$d$日之前在$l$日交付。来源方假设
\[
D_{o,d,e,l}\ \text{and}\ D_{o',d',e',l'}
\ \text{are independent for }e\ne e'.
\]
各周期中都会实现需求量的交付。从发货人到收货人的运输属于城市内部的操作，与这个模型无关。截止日期并不是硬性规定：如果延迟交付，将会受到处罚；而那些可能逾期交付的货物则会被外包处理。

系统中每种商品的总停留时间不得超过$T$。在新的商品出现之前，必须先由运输商或第三方将现有商品交付完毕。托盘只有通过交付到最终目的地或外包给第三方后，才能离开该系统。

\subsection*{状态与决策定义}
对于决策阶段
\[
t\in\mathbb{Z}_{\geq0},\qquad [t]:=t\bmod T,
\]
该状态是
\[
S_t=([t],R_t),\qquad
R_t=(R_{it}^{dl})_{l\in\mathcal{T},\,i,d\in\mathcal{L},\,i\ne d}.
\]
$R_{it}^{dl}$统计终端处的托盘数量；$i$统计纪元开始时的数量；$t$必须于$d$日交付到终端的托盘数量；$l$。库存不区分来源或可用时间。$[t]$日表示可用的拖车容量以及重复需求模式中的位置。

路由决策向量是
\[
x_t=(x_{ijt}^{dl})_{l\in\mathcal{T},\,i,j,d\in\mathcal{L}}.
\]
对于$((i,[t]),(j,[t+1]))\in\mathcal{A}$，$x_{ijt}^{dl}$指的是那些目的地为$d$，交货日期为$l$的托盘：其中，$i\ne j$表示运输过程，而$i=j$则表示在终端站$i$的暂存。路由选择是唯一的决策；外包只是过渡过程中的一部分，并非独立的决策。

\paragraph{印刷索引范围。}
状态向量仅适用于$i\ne d$，而(1)中的第一行则适用于所有$i,d\in\mathcal{L}$，包括$i=d$。所打印的范围被保留下来；这里没有为$R_{dt}^{dl}$提供任何值来修复这种不匹配。同样，决策向量和整体性行也适用于所有$i,j,d\in\mathcal{L}$，而运动解释和流量总和则指的是网络弧线。没有额外的非弧线零约束条件被施加。

\subsection*{确定性决策后的状态转移与随机性状态转换}
在做出可行的决策$x_t\in\mathcal{X}(S_t)$之后，系统会先确定性地从$R_t$转移到决策后的库存状态$\bar R_t=(\bar R_{it}^{dl})$，之后才会处理新的货物。
\begin{equation}
\bar R_{i,t}^{dl}=
\begin{cases}
\displaystyle\sum_{j\in\mathcal{L}:((j,[t]),(i,[t+1]))\in\mathcal{A}}x_{jit}^{dl},
&\begin{gathered}
\forall d\in\mathcal{L},\ i\in\mathcal{L},\ l\in\mathcal{T},\ i\ne d,\\[-2pt]
\Delta^{out}(t+1,i,d,p)=0,
\end{gathered}\\
0,&\text{otherwise}.
\end{cases}
\tag{2}
\end{equation}
固定参数$p$控制延迟容忍度。函数$\Delta^{out}(t+1,i,d,p)$在托盘将被外包的情况下返回1，这意味着其最早到达时间为$p$天或更多，且晚于截止日期$l$，这是基于从时间$t+1$开始，在$\mathcal{G}$中从$i$到$d$的最短路径；否则返回0。在循环网络模型中，来源方表示该函数首先根据逗留时间假设判断截止日期$l$是在过去还是未来，然后判断外包是否不可避免。一旦确定这一点，就会立即评估外包情况。对于在时刻$t$从$i$发送到$j$的托盘，外包惩罚会在$t$进行评估，并且在时刻$t+1$之前将托盘移除。方程（2）也适用于运送到最终目的地的托盘。

\paragraph{印刷外包表示法和边界。}
虽然其定义取决于特定日期$l$，但$\Delta^{out}(t+1,i,d,p)$中的打印参数列表却省略了$l$；不过，该列表本身仍然被保留了下来。 前面的文字说托盘可延迟“至多 $p$ 天”才算过期，但式 (2) 后的定义写的是“$p$ 天或更久”，且状态转移说明中的例子在延迟 $p$ 天时即外包。 这里的操作内容是指(2)中的定义，包括$p$处的等式；没有进行任何阈值修复。

随机到达转移是
\begin{equation}
R_{i,t+1}^{dl}=\bar R_{i,t}^{dl}+D_{i,d,[t+1],l}
\qquad\forall d\in\mathcal{L},\ l\in\mathcal{T},\ i\ne d.
\tag{3}
\end{equation}
这里$i$是库存定义中的终端位置；显示的量化因子保持原样。下一个完整状态由$t+1$时刻的相同状态定义决定，周期为$[t+1]$天。

\subsection*{单期奖励}
对于状态$S_t$和可行的路由$x_t$，这三个打印出来的组件是：
\begin{align*}
C_t^1(S_t,x_t)
&=\beta\sum_{d,l,i,j}x_{ijt}^{dl}\,1_{\{j=d\}},\\
C_t^2(S_t,x_t)
&=-c^{late}\sum_{d,l,i,j}\Lambda(t,l,p)x_{ijt}^{dl}
\,1_{\{j=d\,\wedge\,\Delta^{late}(t,l,p)=1\}},\\
C_t^3(S_t,x_t)
&=-c^{out}\sum_{d,l,i,j}\delta(j,d)x_{ijt}^{dl}
\,1_{\{j\ne d\,\wedge\,\Delta^{out}(t+1,j,d,p)=1\}},\\
C_t(S_t,x_t)
&=C_t^1(S_t,x_t)+C_t^2(S_t,x_t)+C_t^3(S_t,x_t).
\end{align*}
这些求和结果保留了原文中的简写索引$d,l,i,j$，这些索引对应于终端索引$d,i,j\in\mathcal{L}$和到期日$l\in\mathcal{T}$。而因子$1_{\{\cdot\}}$则代表了它们所代表的特定条件。

\begin{itemize}
\item $\beta$指的是承运人将货物交付到最终目的地时所获得的收入。外包运输服务并不产生这种收入。
\item $c^{late}$表示每个托盘在到达最终目的地$d$时迟到的惩罚金额。函数$\Delta^{late}(t,l,p)$在决策时期$t$结束时，如果托盘会迟到，则返回1；否则返回0。函数$\Lambda(t,l,p)$则返回托盘迟到的天数，如果满足$j=d$和$\Delta^{late}(t,l,p)=1$的条件，则返回该天数；否则返回0。这些函数以论文中使用的抽象形式保留下来。
\item $c^{out}$是每托盘外包的惩罚费用。来源将$\delta(j,d)$描述为从终端$j$到目的地$d$在$\mathcal{G}^{flat}$中的最短路径，并将其作为剩余的行程距离：外包惩罚费用与该距离成正比。没有额外的边长约定。
\end{itemize}
固定参数$\beta$、$c^{late}$和$c^{out}$对所有商品都是相同的。仅在计算实验中使用的数值和需求分布并不适用于这一公式。

\subsection*{关联的最优性方程}
该论文给出了针对这一相同的折现马尔可夫决策过程所适用的贝尔曼特征描述以及最优决策规则：
\begin{equation}
V^*(S_t)=\max_{x_t\in\mathcal{X}(S_t)}
\left\{C_t(S_t,x_t)+\gamma\mathbb{E}\left[
\left.V^*(S_{t+1})\right|S_t\right]\right\},
\qquad\forall S_t.
\tag{4}
\end{equation}
\begin{equation}
X^{\pi^*}(S_t)=\arg\max_{x_t\in\mathcal{X}(S_t)}
\left\{C_t(S_t,x_t)+\gamma\mathbb{E}\left[
\left.V^*(S_{t+1})\right|S_t\right]\right\},
\qquad\forall S_t.
\tag{5}
\end{equation}
$V^*$是从某个状态开始，未来预期获得的总奖励的最佳值。下一个状态是通过(2)--(3)公式在选定的动作下获得的。而(4)--(5)中的条件设定则保持原样。

\paragraph{来源范围。}
本文转录 Baubaid、Boland 与 Savelsbergh 的论文 \emph{零担运输承运人的动态货运路径问题} 中第 3.1 至 3.2.5 节的模型，日期为 2020 年 11 月 5 日，PDF 实际第 7 至 12 页；可用时以所提供的页面图像为准。 所引用的关于设计 2-alt 负载方案的配套论文并未提供，因此在此无法验证其准确性；该论文中的设计方案并未被用作固定路由集 $\mathcal{Y}$的替代方案。 没有提供官方补充材料。 ADP近似方法、扩展的整数编程子问题、聚合方案以及实验性实例限制等，都超出了这个选定模型的范畴。

\end{document}
